\documentclass[12pt,a4paper]{article}
\usepackage[utf8]{inputenc}
\usepackage[T1]{fontenc}
\usepackage[spanish]{babel}
\usepackage{amsmath,amssymb,amsfonts}
\usepackage{graphicx}
\usepackage{geometry}
\usepackage{mathpazo}
\usepackage{float}
\usepackage{hyperref}

\geometry{margin=2.5cm}
\linespread{1.05}
\setlength{\parskip}{0.5em}

\title{Título del Capítulo}
\author{Autor}
\date{\today}

\begin{document}

\maketitle

\section{Introducción}

Este documento sirve como plantilla base para redactar capítulos dentro del proyecto \texttt{Navier\_Stokes\_Analitico\_Numerico}. Aquí podés comenzar a escribir el contenido matemático o físico en español técnico.

\section{Sección de ejemplo}

Un ejemplo de ecuación alineada:

\begin{align}
\frac{\partial \rho}{\partial t} + \nabla \cdot (\rho \mathbf{v}) &= 0, \\
\frac{\partial (\rho \mathbf{v})}{\partial t} + \nabla \cdot (\rho \mathbf{v} \otimes \mathbf{v}) &= -\nabla p + \mu \Delta \mathbf{v}.
\end{align}

\section{Figuras y referencias}

\begin{figure}[H]
\centering
\includegraphics[width=0.6\textwidth]{ruta/figura.pdf}
\caption{Ejemplo de figura}
\label{fig:ejemplo}
\end{figure}

Podés referenciar la Figura~\ref{fig:ejemplo} o ecuaciones como la (1).

\section{Notas finales}

Recordá que al guardar, se compila automáticamente con \texttt{latexmk} si usás VS Code y LaTeX Workshop.

\end{document}